\section{Aufbau}
%\subsection{Winkelverteilung}
%abbildung von aufbau einfügen
Zur Messung der Winkelverteilung kosmischer Myonen wird ein Detektor mit 24 planaren Szintillatoren und einem zentralen zylindrischen Szintillator verwendet. Die äußeren Detektoren sind gleichmäßig auf einem 30 cm-Raster verteilt (siehe Abbildung ???). Jeder Scan entspricht einem festen Azimut von 15°. Der zentrale Szintillator dient als Datendetektor und zeigt zusammen mit den beiden äußeren Detektoren die geometrische Position über die Zeit an.

Die Ereignisse wurden mittels Drei-Koinzidenz-Darstellung erfasst. Dazu werden die analogen Ausgangssignale aller Szintillatoren zunächst über Vorverstärker an Differenzwandler geleitet, die sie in digitale Logikimpulse mit geringer Amplitude umwandeln. Die Differenz zwischen den Signalen der beiden äußeren Detektoren und denen des zentralen Detektors (Z25) bildet die Koinzidenzeinheit. Dieses Verfahren liefert nur dann ein korrektes Ergebnis, wenn das Logiksignal innerhalb eines vorgegebenen Zeitintervalls gleichzeitig an allen drei Eingängen anliegt. Die Koinzidenzeinheit entspricht Myonen, die das Gerät in zufälliger Richtung passieren und denen ein bestimmter Azimutwinkel zugeordnet werden kann.

Durch die Untersuchung zweier verschiedener Detektorgruppen lässt sich die Flugbahn der Myonen in der oberen Hemisphäre berechnen. Das Produkt der Zählwerte über ein 15°-Intervall ergibt die Verteilung der Höhenstrahlung. 

%\subsection{Myonenlebensdauer}
Im zweiten Teil des Experiments wird die Lebensdauer des Myons ermittelt. Hierfür wurde der in Abbildung ??? dargestellte Versuchsaufbau verwendet. Eine 5 cm dicke Bleiplatte schirmt den schmalen Photonenstrahl ab und lenkt ihn auf das Substrat.

Wenn ein einfallendes Photon auf zwei übereinanderliegende Detektoren trifft, emittiert es gleichzeitig in beiden ein elektrisches Signal. Dieses Signal dient als Startsignal für die Lebensdauermessungen. Der untere Szintillationsdetektor (Z2) ist massereich genug, um einige Myonen einzufangen und zu neutralisieren. Der bei der Kollision erzeugte elektrische Impuls erzeugt in Z2 ein zweites Signal, das als Stoppsignal aufgezeichnet wird.

Die Zeitdifferenz zwischen Start- und Stoppsignal wird manuell erfasst und durch eine von zehn $1\mu s$-Taktperioden dividiert. Die Zählrate des Detektors folgt dem erwarteten Trend, wodurch die Myonlebensdauer bestimmt werden kann.

Da das Gerät das Stoppsignal eines einfallenden Myons nicht von dem seines Zerfalls unterscheiden kann, wird für jeden Startimpuls ohne weitere Verarbeitung ein Stoppsignal erzeugt. Um dieses Problem zu beheben, wird das Startsignal mithilfe eines Verzögerungskabels um 100 ns gegenüber dem Stoppsignal verzögert. Dadurch werden die Stoppsignale der einfallenden Myonen zufällig. Diese Maßnahme bewirkt automatisch eine lineare Variation in der gemessenen Zeitverteilung, beeinflusst aber weder das exponentielle Verhalten noch die durch Anpassung ermittelten Lebensdauern.
